\section*{2026 年 4 月 20 日作业}
\phantomsection
\addcontentsline{toc}{section}{2026 年 4 月 20 日作业}
\noindent\textbf{截止时间:2026 年 4 月 27 日 13:30}
\setcounter{problem}{0}
\setcounter{sollemma}{0}

\begin{problem}
已知 $n$ 是正整数, $A$ 是 $n$ 阶 Hermite 矩阵.证明: $A$ 是半正定矩阵的充要条件是 $A$ 的所有主子式均非负.
\begin{solution}
    证明:\textbf{充分性($\Rightarrow$)}:

    \textbf{必要性($\Leftarrow$)}:必要性显然.任取 $k \in $
\end{solution}
\end{problem}

\begin{problem}
已知 $n$ 是正整数,若 $A,B,A-B$ 都是 $n$ 阶 Hermite 正定矩阵.证明:$B^{-1}-A^{-1}$ 也是 Hermite 正定矩阵.
\begin{solution}
    证明:
\end{solution}
\end{problem}

\begin{problem}
已知 $n$ 是正整数, $A,B$ 是 $n$ 阶 Hermite 矩阵, $B$ 正定.证明:存在 $n$ 阶非奇异矩阵 $X$ 使得 $X^*AX$ 和 $X^*BX$ 都是对角矩阵.
\begin{solution}
\end{solution}
\end{problem}

\begin{problem}
已知 $n$ 是正整数,若 $A,B$ 都是 $n$ 阶 Hermite 矩阵.证明: $\operatorname{In}_{+}(A+B)\le \operatorname{In}_{+}(A)+\operatorname{In}_{+}(B)$.
\begin{solution}

\end{solution}
\end{problem}

\begin{problem}
已知 $n$ 是正整数,可定义 $n$ 阶 Pascal 矩阵:
\[
    P_n=
    \begin{bmatrix}
        1 & 1 & 1 & 1 & \cdots \\
        1 & 2 & 3 & 4 & \cdots \\
        1 & 3 & 6 & 10 & \cdots \\
        1 & 4 & 10 & 20 & \cdots \\
        \vdots & \vdots & \vdots & \vdots & \ddots
    \end{bmatrix}
    \in\R^{n\times n},
\]
其 $(i,j)$ 分量为
\[
    \binom{i+j-2}{i-1}
    =
    \frac{(i+j-2)!}{(i-1)!(j-1)!}.
\]
证明 $P_n$ 是正定矩阵,并求其 Cholesky 分解.
\begin{solution}

\end{solution}
\end{problem}

\begin{problem}[选做]
已知 $m,n$ 是正整数,若 $A\in\C^{n\times n}$ 是 Hermite 半正定矩阵.证明:存在唯一的 Hermite 半正定矩阵 $X$ 使得 $X^m=A$.
\begin{solution}

\end{solution}
\end{problem}

\begin{problem}[选做]
证明 Schur 乘积定理:已知 $n$ 是正整数,若 $A,B\in\C^{n\times n}$ 都是 Hermite 正定矩阵,那么 $C=A\circ B$ 也是 Hermite 正定矩阵.这里 $\circ$ 表示矩阵的 Hadamard 乘积.
\begin{solution}

\end{solution}
\end{problem}

\begin{problem}[选做]
设 $n$ 是正整数,$x,y\in\R^n\setminus\{0\}$.证明:$y^Tx>0$ 的充要条件是存在对称正定矩阵 $A\in\R^{n\times n}$ 使得 $y=Ax$.
\begin{solution}

\end{solution}
\end{problem}
